\section{KẾT LUẬN}
Trong bài tập lớn này, nhóm đã khảo sát các phương pháp tăng cường ảnh, khử nhiễu và thứ tự kết hợp khi nhận dạng 43 lớp biển báo GTSRB. Nhóm tạo ảnh thiếu sáng bằng Gamma kết hợp Poisson và Gaussian ở ba mức. Các phép xử lý sử dụng cùng ảnh đầu vào và ba checkpoint ResNet-18 đã giữ cố định. Nhóm lựa chọn tham số và danh sách phương pháp đánh giá toàn bộ test bằng validation trước khi chạy test.

Trong nhóm tăng sáng, diffusion được chọn theo quy tắc đánh giá validation. Khi xét thêm khử nhiễu và kết hợp, chuỗi Gaussian $5\times5$, $\sigma=0,8$ rồi Gamma $\alpha=0,8$ đạt Macro-F1 validation 77,77\% và được chọn. Trên 12.630 ảnh test, chuỗi đạt Macro-F1 lần lượt là 96,70\%, 85,54\% và 71,56\% ở ba mức suy giảm. So với ảnh tối giữ nguyên, mức tăng tương ứng là 6,25, 9,47 và 13,19 điểm phần trăm. Dù có cải thiện, ảnh thiếu sáng nặng vẫn là trường hợp khó đối với mô hình.

Qua kết quả, nhóm nhận thấy PSNR và SSIM chưa thể thay thế đánh giá nhận dạng. Ở mức nặng, Gaussian đơn lẻ có cả PSNR và SSIM thấp hơn diffusion nhưng F1 cao hơn. Trong thí nghiệm ablation, MSR đạt 98,36\% khi chỉ có Gamma làm tối, nhưng còn 57,84\% khi ảnh có cả hai loại nhiễu. Những kết quả này cho thấy ảnh sáng hơn hoặc có một chỉ số ảnh tốt hơn chưa đủ để kết luận các chi tiết phân biệt lớp đã được giữ lại.

Ảnh hưởng của thứ tự còn phụ thuộc bộ lọc, phép tăng sáng và cấu hình cụ thể. Trên toàn bộ test, cặp Gaussian--CLAHE có chênh lệch trung bình nhỏ theo hướng CLAHE rồi Gaussian tốt hơn, nhưng một số lượt lại cho kết quả ngược lại. Các khoảng bootstrap 95\% chỉ lấy mẫu ảnh nằm trên 0, còn các khoảng lấy mẫu thêm checkpoint đều chứa 0. Vì vậy, nhóm chưa có đủ bằng chứng để khẳng định ưu thế này ổn định khi huấn luyện lại. Nhóm cũng chưa xác định được một thứ tự tốt nhất cho mọi phương pháp. Ở mức vừa, chuỗi được chọn sửa đúng trung bình 1472,67 ảnh nhưng làm sai 29,22 ảnh mỗi lượt. Do đó, khi đánh giá lợi ích, nhóm cần xem cả những lỗi mới phát sinh.

Đóng góp của nhóm là triển khai, kiểm tra và so sánh các phép xử lý trong điều kiện có kiểm soát, chưa phải đề xuất thuật toán mới. Các kết quả giúp nhóm lựa chọn tiền xử lý theo khả năng nhận dạng thay vì chỉ nhìn độ sáng đầu ra. Tuy nhiên, kết luận mới áp dụng cho ROI GTSRB với suy giảm mô phỏng, một kiến trúc CNN và lưới tham số hữu hạn. Nhóm chưa chứng minh hiệu quả trên camera ban đêm thực, video hoặc hệ thống cảnh báo trong xe. Đây là những phần cần bổ sung nếu phát triển đề tài theo hướng ứng dụng.
\begin{center}\vspace{6mm}---------- HẾT ----------\end{center}
